\scanpage{98}
% archive foliation: 97
un op. \unsure{1-diff.} (rel. à $X/S$)
\[
D:\underline{\omega}\longrightarrow L,
\]
\deleted{satisfaisant}
et la restriction de $\gamma$ à $\underline{\omega}\otimes\underline{\omega}=\underline{\omega}^2\subset(P^1_{X/S}\omega)$
\[
\gamma_0:\underline{\omega}^2\unsure{\text{---}|\longrightarrow}L\qquad\text{i.e. }\gamma_0\in\Gamma(L\otimes\omega^{-2})
\]
est aussi définie par
\[
D(f\omega)-fD(\omega)=\gamma_0(\omega\cdot df)
\]
($\omega\in\Gamma\underline{\omega}$, $f\in\Gamma\underline{\mathcal{O}}_X$). Les pts où la section $\sigma'$ déduite de $\gamma$ (ou $D$) rencontre la section $\sigma$ canonique sont ceux où $\gamma_0$ s'annule, i.e. les \uline{pts de dégénérescence} de l'op. diff. $D$. On notera qu'aux pts de non dég., il est automatique que $\gamma$ est épi -- que ce soit épi en tt pt signifie aussi que $L$ est engendré comme module par les $D(\omega)$ ($\omega$ sections locales de $\underline{\omega}$). Si ce n'est pas vrai, alors (si $X$ courbe alg. sur corps de base $k$) il y a un sous-Module $L'\subset L$ canonique par lequel se factorise $D$ (le plus petit --) qui est inv. et pour lequel la condition est vérifiée --

On voit donc qu'à tt tel épi est associée, \uline{sur l'ouvert des pts où elle est non dégénérée}, une connexion projective \deleted{\unsure{qui}} unique qui la rende horizontale. \deleted{La question si $\exists$ connexion}
